\documentclass[12pt,a4paper]{article}
\usepackage[margin=2.2cm]{geometry}
\usepackage{xcolor}
\usepackage{setspace}
\usepackage[colorlinks=true,urlcolor=blue!60!black]{hyperref}
\usepackage{microtype}

\newcommand{\slide}[3]{%
  \vspace{10pt}
  \noindent{\large\bfseries Slide #1 --- #2}\hfill{\color{gray}\small #3}\par
  \vspace{2pt}\noindent\rule{\textwidth}{0.4pt}\par\vspace{4pt}}

\setstretch{1.15}
\setlength{\parindent}{0pt}

\begin{document}

\begin{center}
{\LARGE\bfseries Seminar Speaker Script}\\[4pt]
{\large Sparse Recovery and Compressed Sensing: A Model-Based Deep Learning Approach}\\[4pt]
Model-Based Deep Learning (361.2.2320) --- Recorded seminar, $\sim$15 minutes\\[2pt]
{\color{gray} Etay Baron \& Etai Wigman. The same text is embedded in the PPTX as per-slide speaker notes.}
\end{center}

\slide{1}{Title}{0:00--0:30}
Hello everyone. We are Etay Baron and Etai Wigman, and this is our final project for the Model-Based Deep Learning course. Our topic is sparse recovery and compressed sensing, based on the NeurIPS 2021 paper ``Hyperparameter Tuning is All You Need for LISTA.'' The core question that guides this talk is simple: how far can model-based structure take us --- and when do we need data-driven flexibility?

\slide{2}{Seminar Roadmap}{0:30--1:10}
Here is the plan for the next fifteen minutes. We start with the background: the sparse recovery problem, the LASSO formulation, and the classical ISTA and FISTA algorithms. We then move to deep unrolling --- LISTA, ALISTA, and HyperLISTA. The central part is our own study: we ask which components of ISTA are actually worth learning, and we compare the training methods for unfolded optimizers that were discussed in class. Finally, we stress-test everything on real images --- MNIST and Fashion-MNIST --- under compressed sensing. Our synthetic setup uses 250 measurements, signal dimension 500, and sparsity 50.

\slide{3}{Motivation \& Problem Formulation}{1:10--2:10}
Let us define the problem. We want to recover a sparse vector $x^*$ of dimension $n$, but we only observe $m$ linear measurements, with $m$ smaller than $n$, possibly with noise. This inverse problem is underdetermined --- infinitely many signals explain the same measurements. The sparsity assumption is what makes it solvable, and this situation appears everywhere: MRI acceleration, radar, hyperspectral imaging. The standard formulation is the LASSO: a trade-off between fidelity to the measurements and an $\ell_1$ norm that promotes sparsity. Our metrics are NMSE in decibels for synthetic signals, plus PSNR and SSIM for images.

\slide{4}{Classical Baselines: ISTA and FISTA}{2:10--3:10}
ISTA is proximal gradient descent on the LASSO objective. Each iteration takes a gradient step on the data-fidelity term, and then applies soft-thresholding, which shrinks small entries to zero. FISTA adds Nesterov momentum and improves the convergence rate from order $1/k$ to $1/k^2$. But here is the key point: after only sixteen iterations, ISTA reaches just minus 5.3~dB, and FISTA minus 11. Reaching minus 20~dB takes hundreds of iterations. So the question is: how do we make sixteen steps really count?

\slide{5}{Deep Unrolling: From Iterations to Layers}{3:10--4:20}
The idea, due to Gregor and LeCun in 2010, is to unfold the $K$ iterations of ISTA into a network with $K$ layers, and to turn the algorithmic constants into learnable parameters. On the left you see classical ISTA, with a fixed step size and a fixed threshold. On the right is LISTA: two weight matrices and a separate threshold per layer, all trained end-to-end with backpropagation. This is the model-based deep learning principle: keep the structure --- a gradient step followed by soft-thresholding --- and learn only what improves reconstruction at a finite depth. In our setting this amounts to about six million parameters.

\slide{6}{The MBDL Ladder: LISTA $\rightarrow$ ALISTA $\rightarrow$ HyperLISTA}{4:20--5:50}
Now the MBDL ladder --- more structure, fewer parameters. LISTA learns everything: six million parameters, minus 20~dB. ALISTA shows that the weight matrix does not need to be learned at all: it can be computed analytically from $A$, leaving only per-layer step sizes and thresholds --- 32 parameters, and minus 30~dB. HyperLISTA takes this to the extreme: only three scalars. The threshold, the momentum, and the support selection are computed adaptively from the residual at every layer, using the formulas on the slide. And notice --- there is no backpropagation at all: the three scalars are found by a simple grid search. The result: minus 61~dB, the best of all. More structure, fewer parameters, better performance --- this is the MBDL principle in action.

\slide{7}{Our Design Space Study: What Should We Learn?}{5:50--6:50}
This leads to our research question: between six million parameters and three, what is the minimal set that is actually worth learning? The generalized ISTA update exposes two free scalar sequences: the step size $\gamma$ and the threshold $\theta$. We designed three models. ThresholdISTA learns only the thresholds. StepISTA learns only the step sizes, with the threshold tied to them. And StepThresholdISTA learns both, independently --- 32 scalars in total. All of them are initialized to the exact ISTA values, so before training, every model is identical to ISTA.

\slide{8}{Our Results: Scalar Models}{6:50--8:00}
The results are quite striking. Learning thresholds alone barely helps: minus 7~dB. But learning the step-size schedule alone --- just sixteen scalars --- jumps to minus 20~dB and matches the full LISTA with six million parameters. Decoupling the threshold from the step adds another 3.6~dB, essentially for free. The conclusion: the step-size schedule is the critical component to learn at finite depth, and learning full matrices is largely redundant when the sensing matrix is fixed.

\slide{9}{Training Methods for Unfolded Optimizers}{8:00--9:10}
Following the training strategies discussed in class, we compared them on the exact same architecture. L1 is standard end-to-end training, with a loss on the final layer only. L2 is deep supervision: a weighted loss on every intermediate layer. L3 is greedy layer-wise training, where each layer is trained while the previous ones are frozen. We also tested weight tying --- the same layer repeated $K$ times, like a recurrent network. The interesting finding: tied LISTA, with sixteen times fewer parameters, comes very close to the untied version. Weight tying acts as a built-in regularizer.

\slide{10}{Synthetic Sparse Recovery Results}{9:10--10:10}
Here is the full picture on the synthetic problem: NMSE as a function of depth. The training method matters --- L2 is the best, and L3 gets stuck around layer eight, because later layers have no incentive to improve on what earlier layers already solved. But the big hierarchy remains: ALISTA beats every LISTA variant with 32 parameters, and HyperLISTA beats everyone with three. When the model assumptions hold, structure beats capacity.

\slide{11}{Design-Space Map}{10:10--11:00}
This map summarizes the synthetic part. The x-axis is the number of parameters on a logarithmic scale; the y-axis is the reconstruction error; the bottom-left corner is the ideal. Our StepThresholdISTA, with 32 scalars, sits in an excellent spot and beats LISTA with 187,000 times fewer parameters. HyperLISTA breaks the scale entirely. The message: in problems with known structure, choosing what to learn matters more than raw capacity.

\slide{12}{Real Image CS: FashionMNIST Pixel Domain}{11:00--12:10}
Now the real test --- actual images. We moved to pixel-domain compressed sensing on Fashion-MNIST: the image is flattened into a 784-dimensional vector, and we measure only a quarter of that dimension. No transform is needed, because these images have a black background and are naturally sparse in pixels. And here the ranking flips completely. All the structured methods collapse --- ALISTA and HyperLISTA do even worse than plain ISTA --- while LISTA with deep supervision reaches an SSIM of 0.82. The reason: the assumption of i.i.d.\ Gaussian sparsity is simply wrong for bounded, spatially-correlated pixels. A wrong prior is worse than no prior.

\slide{13}{Robustness Across Datasets and Measurement Ratios}{12:10--13:10}
We tested this systematically over two datasets and two measurement ratios; the heatmap shows SSIM. On MNIST with half the measurements, every method does fine, because MNIST is genuinely sparse --- about 78 percent of its pixels are zero. As the task gets harder --- fewer measurements, or more textured data --- the structured methods degrade sharply, while LISTA-L2 stays above 0.82 in every scenario. The sparse prior breaks down gradually as the data moves away from the assumptions.

\slide{14}{Visual Stress Tests}{13:10--14:00}
The visual comparison makes the numbers concrete. MNIST is on the left, Fashion-MNIST on the right, both with a quarter of the measurements. The LISTA-L2 row is almost identical to the originals in both cases. ALISTA and HyperLISTA literally erase real signal --- a threshold calibrated for sparse Gaussian residuals deletes the image itself. This is what a wrong prior looks like.

\slide{15}{Conclusions}{14:00--14:50}
To summarize, four lessons. One: when the assumptions hold, structure wins --- three scalars beat six million parameters by 41~dB. Two: the step-size schedule is the most valuable thing to learn. Three: weight tying is implicit regularization. Four: a wrong prior actively hurts --- on real images, the flexible LISTA wins. As future work, we would like to test these models in the DCT or wavelet domain, where images are much closer to the Gaussian-sparse assumption.

\slide{16}{Code \& References}{14:50--15:00}
All of our code --- the four notebooks and the source modules --- is available in a public GitHub repository, and these are the main references we built on. Thank you very much for listening.

\end{document}
